\section{Summary}
This thesis investigated the influence of AGN on colour-colour diagrams used for redshift determination, galaxy classification, and AGN identification. To achieve this, we selected representative models of Type 1 and Type 2 AGN, reflecting the primary AGN categories relevant to this study. These models were then integrated with a suite of state-of-the-art galaxy templates to generate a composite set of spectral energy distributions (SEDs) with varying AGN contributions. By generating synthetic photometry across relevant filter bands, we analysed how the choice of the AGN model affected the composite galaxy's position within various colour-colour spaces. This approach allowed us to assess the potential for AGN to contaminate the inherent colours of a galaxy, depending on the contribution amount. \\\\ 
Using these theoretical models, we found that Type 1 AGN significantly impact the UVJ, ugr, and IRAC colour-colour diagrams, even at very low AGN contributions. Further investigation revealed that even a minor increase in Type 1 AGN contribution could lead to an overestimation of sources appearing as star-forming galaxies in the UVJ diagram. This is because quiescent galaxies were shown to move through the colour space to the star-forming region with increasing AGN contribution. Similarly, we found that small increases in AGN contribution can also cause redshift estimation using the ugr diagram to fail, causing sources to appear in lower redshift regions.  In the IRAC diagram, our type 2 models behaved as expected and were correctly selected by the Lacy Wedge, while Type 1 AGN were removed from the selection region. \\\\
To ensure our modelling approach was applicable to observable galaxies at high redshifts, we attempted to replicate these results with galaxy SEDs constructed using EAZY from ZFOURGE photometry, the result of which showed very similar trends. Following this, the theoretical results were compared against a collection of SEDs fitted to ZFOURGE photometry, with their AGN component subsequently removed using CIGALE. This decomposed observational data showed that, in general, galaxies without their AGN component appear more quiescent in the UVJ diagram – a result consistent with the theoretical modelling. Observational results for the ugr diagram could not be reproduced due to inconsistencies in the photometry, an issue that will be addressed in future work. This research highlights the potential for Type 1 AGN to distort galaxy classification and redshift estimation, underscoring the need for robust diagnostics, particularly in the context of high-redshift galaxy studies.
\newpage
\section{Future Directions}
% Extension to other colour spaces and custom colour-colour spaces
While the work was limited to the UVJ, ugr, and IRAC colour spaces, other useful spaces could be investigated. For instance, examining the gri colour space could reveal how Type 1 AGN may impact the identification of high-redshift sources using the g-dropout technique \citep{giavalisco_lyman-break_2002}. Similarly, further work could explore the effect of both Type 1 and Type 2 AGN on WISE colour space \citep{wright_wide-field_2010}. This would be useful as WISE colours are versatile for identifying AGN, separating and classifying types of astronomical objects \citep{krakowski_machine-learning_2016}, and studying star formation \citep{caccianiga_wise_2015}.
Furthermore, employing synthetic James Webb Space Telescope (JWST; \citealt{gardner_james_2006}) could yield crucial insights into how AGN influence JWST colour spaces. This is particularly relevant considering the wealth of high-quality data from JWST, which has facilitated galaxy evolution studies of the most distant galaxies. Modelling AGN effects on these filters will reveal any significant impact on the colour space, proving valuable for future studies employing JWST photometry. To achieve a comprehensive and robust understanding of the AGN population, future research must incorporate intermediate-type AGN, as they represent a substantial fraction of observed objects and are crucial for explaining the diverse classes within this category. Building upon these findings, future investigations could explore how machine learning algorithms might dynamically design selection regions to mitigate the effects of AGN on the colour space.  Other avenues include leveraging these results to remove AGN using advanced classification algorithms such as Support Vector Machines (SVM; \citealt{hearst_support_1998}) or XGBoost \citep{chen_xgboost_2016}.

% Investigating the decomposed seds degeneracy and other ways of observational confirmation